\chapter*{How to Read This Report}
\addcontentsline{toc}{chapter}{How to Read This Report}
\markboth{How to Read This Report}{}

This report explains the VidyaRAG repository end to end: what it does, how every stage works, how it was measured, how it ships, where it is weak, and how to talk about it in an interview. It was written from the code, not from the documentation. Where the two disagreed, the code won and the disagreement is recorded in Chapter~\ref{ch:lessons}.

\paragraph{What it describes.} Version~1.3.0 of the repository, after the measurement fixes of pull request~\#20. It began as an audit of commit \code{fa5164b} (version 1.2.1, 14 September 2026). That audit found 32 defects, including one that invalidated the project's headline result. The fixes that followed are the reason this report can state corrected numbers as current fact. Chapter~\ref{ch:lessons} tells that story, because it is the most useful thing in the report for an interview.

\paragraph{The kinds of statement.} They are never mixed inside one sentence.
\begin{description}[style=sameline,leftmargin=3.4cm,font=\normalfont\bfseries]
  \item[Verified] Plain text, tables and figures. Checked against the code, a config file, a test, a committed result file or git history. Non-obvious claims carry a grey file reference such as \ev{src/vidyarag/pipeline.py}.
  \item[Measured] Teal boxes, or the tag \tagM. Produced by running the system while preparing this report, on a Windows~11 machine without a GPU, Python 3.11, with the Gemini API on the free tier. These numbers are not in the repository.
  \item[Inferred] Orange boxes, or the tag \tagI. Reasoned from the code but neither stated nor measured.
  \item[Recommended] Purple boxes, or the tag \tagR. Not implemented.
\end{description}
Grey \emph{background} boxes hold general knowledge that is not about this repository. Red \emph{finding} boxes mark a defect or contradiction. When the repository cannot answer a question, the report says \NotDet{} It never guesses.

\paragraph{Rules about numbers.} No number is invented. Each one is quoted from a named committed file, computed from committed files by a stated method, or measured and labelled as such. All evaluation results are single runs over a 58-question set, so Chapter~\ref{ch:evaluation} says which differences exceed the measured run-to-run noise.

\paragraph{Notation.} \code{Monospace} is code, paths and exact strings. A \emph{profile} is one of five committed pipeline configurations: \code{baseline}, \code{rerank}, \code{decompose}, \code{corrective} and \code{guarded}. The demo runs \code{guarded}. \emph{Chunk} and \emph{passage} both mean one indexed piece of text. Costs are US dollars at the list prices written in the code, not money spent: the project runs on a free tier.

\paragraph{Reading paths.}
\begin{center}\small
\begin{tabularx}{\textwidth}{@{}p{4.0cm}L@{}}
\toprule
\textbf{If you have\ldots} & \textbf{Read}\\
\midrule
One hour & Chapters~\ref{ch:overview}, \ref{ch:architecture}, \ref{ch:lessons} and~\ref{ch:interview}.\\
One evening & Chapters~1--3, \ref{ch:journey} (one question end to end), \ref{ch:evaluation} and~\ref{ch:lessons}.\\
The day before an interview & Chapter~\ref{ch:interview}, then Chapters~\ref{ch:lessons} and~\ref{ch:judgement}.\\
A plan to extend it & Chapters~\ref{ch:stack} to~\ref{ch:journey} in order, then Chapter~\ref{ch:shipping}.\\
\bottomrule
\end{tabularx}
\end{center}
